\section{A balanced vacuum and its Lie--chiral comparison}
\label{bv:sec}

The ordinary calculation of Proposition~\ref{prop:sl2} suggests a
cohomology class in degree three, but a comparison of vertex objects
must also preserve translation and every mode product.  An explicit
example is obtained by coupling two universal Virasoro vacua of charges
\(a\) and \(26-a\) to free ghosts.  The charge balance makes the
differential square-zero.  A weight contraction then reduces both the
balanced complex and a smaller cochain vertex algebra to the same four
states.  The resulting maps also compare their Lie--chiral Chevalley
coalgebras, with every word length retained.

\subsection{Universal vacua, ghosts and the differential}

Let \(R\) be a nonzero commutative \(\mathbb Q\)-algebra and
\(a\in R\).  All state spaces in this subsection and the next three
are algebraic direct sums of their integral weight spaces.  Cohomological
degree supplies the parity.  Write
\[
 Y(u,z)=\sum_{n\in\mathbb Z}u_{(n)}z^{-n-1},\qquad
 {:}uv{:}=u_{(-1)}v,\qquad
 [u_\lambda v]=\sum_{n\geq0}\frac{\lambda^n}{n!}u_{(n)}v.
\]
The translation operator \(T\) has cohomological degree zero and
raises weight by one.

The universal Virasoro vacuum vertex algebra \(V_R^a\) has an even
stress tensor \(L\), of degree zero and weight two, satisfying
\[
 [L_\lambda L]=(T+2\lambda)L+\frac a{12}\lambda^3\unit.
\]
In physical mode notation \(L_n=L_{(n+1)}\), the vacuum is killed by
\(L_n\) for \(n\geq-1\).  The universal induced vacuum has PBW basis
\[
 L_{-n_1}\cdots L_{-n_s}\unit,
 \qquad n_1\geq\cdots\geq n_s\geq2.
\]
In particular, its weight-zero space is \(R\unit\), its weight-one
space is zero, and every weight space is finite free over \(R\).
This is the universal vacuum, including at charges where it has a
proper vertex ideal.

The free odd ghost vertex algebra has generators \(b,\gamma\), with
degrees \(-1,1\), weights \(2,-1\), and brackets
\[
 [b_\lambda\gamma]=[\gamma_\lambda b]=\unit,
 \qquad [b_\lambda b]=[\gamma_\lambda\gamma]=0.
\]
Its states are the exterior PBW monomials in \(T^rb,T^r\gamma\),
\(r\geq0\).  Put \(\gamma_r=T^r\gamma\).  The subalgebra generated
by the \(\gamma_r\) is the commutative vertex algebra
\[
 G_R=\Lambda_R(\gamma_r:r\geq0),\qquad
 Y(u,z)v=(e^{zT}u)v,
 \qquad |\gamma_r|=1,\quad \operatorname{wt}\gamma_r=r-1.
\]
It has differential \(d\gamma_r=T^r(\gamma_0\gamma_1)\), extended as
an odd derivation.  Indeed \(d\) commutes with \(T\), and
\(d^2\gamma_0=0\) by the exterior relations, so \(d^2=0\).

Use \(L'\) for the stress tensor of \(V_R^{26-a}\).  On
\[
 K_{R,a}=V_R^a\otimes_R V_R^{26-a}\otimes_R bc
\]
define the ghost stress tensor and the odd zero-mode operator
\begin{equation}\label{bv:eq:charge}
 \begin{split}
 L_{\mathrm{gh}}&=-{:}(Tb)\gamma{:}-2{:}bT\gamma{:},\\
 Q&=\bigl({:}\gamma(L+L'){:}+{:}b\gamma T\gamma{:}\bigr)_{(0)}.
 \end{split}
\end{equation}
The notation \(bc\) denotes the free ghost algebra just specified;
its second field is denoted by \(\gamma\).  A zero mode is a
derivation of every vertex product and commutes with \(T\).

\begin{lemma}[The balanced differential]\label{bv:lem:differential}
The operator \(Q\) has degree one, preserves weight and satisfies
\(Q^2=0\).  Write \((L^{(1)},L^{(2)})=(L,L')\) and
\((a_1,a_2)=(a,26-a)\).  Its values on the generating fields are
\begin{equation}\label{bv:eq:generators}
 \begin{split}
 Q\gamma&=\gamma T\gamma,\qquad
 Qb=L+L'+L_{\mathrm{gh}},\\
 QL^{(i)}&=\gamma TL^{(i)}+2(T\gamma)L^{(i)}+\frac{a_i}{12}T^3\gamma,
 \qquad i=1,2.
 \end{split}
\end{equation}
If \(H_K\) denotes total weight on \(K_{R,a}\), then
\begin{equation}\label{bv:eq:cartan-K}
 Qb_{(1)}+b_{(1)}Q=H_K.
\end{equation}
\end{lemma}

\begin{proof}
Single contractions in \eqref{bv:eq:charge} give
\eqref{bv:eq:generators}.  The ghost stress tensor has charge \(-26\):
the four complete double contractions in
\(L_{\mathrm{gh}}(z)L_{\mathrm{gh}}(w)\) contribute
\(-1,-4,-4,-4\) to the fourth pole, whose sum is \(-13\).
The single contractions give weights two and minus one on \(b\)
and \(\gamma\), respectively.

The scalar terms in the differential of the ghost stress tensor must
also be retained.  Wick expansion gives
\[
 \begin{aligned}
 {:}(TL_{\mathrm{gh}})\gamma{:}
   &=3{:}(Tb)\gamma T\gamma{:}
       +2{:}b\gamma T^2\gamma{:}+\frac56T^3\gamma,\\
 {:}L_{\mathrm{gh}}T\gamma{:}
   &=-{:}(Tb)\gamma T\gamma{:}+\frac23T^3\gamma.
 \end{aligned}
\]
For example, before expanding the uncontracted field, the contraction
part of \(L_{\mathrm{gh}}(z)\gamma(w)\) is
\[
 -\frac{\gamma(z)}{(z-w)^2}+\frac{2T\gamma(z)}{z-w}.
\]
Differentiation in \(z\) and extraction of the regular coefficient
give \((1/3-3/2+2)T^3\gamma=5T^3\gamma/6\).
Differentiation in \(w\) gives
\((-1/3+1)T^3\gamma=2T^3\gamma/3\).
Substituting the two identities into \(QL_{\mathrm{gh}}\) cancels
the cubic ghost terms and yields
\[
 QL_{\mathrm{gh}}
 =-\gamma T(L+L')-2(T\gamma)(L+L')-\frac{26}{12}T^3\gamma.
\]
Thus \(Q^2b=0\).  Substitution in
\eqref{bv:eq:generators} gives \(Q^2\gamma=Q^2L=Q^2L'=0\): the
coefficients of \(\gamma(T\gamma)TL^{(i)}\),
\(\gamma T^2\gamma L^{(i)}\), \(\gamma T^4\gamma\) and
\((T\gamma)T^3\gamma\) cancel separately.  These fields generate
\(K_{R,a}\), so the square of the derivation vanishes everywhere.
Finally, the coefficient of mode one in
\([Q,Y(b,z)]_+=Y(Qb,z)\) is \eqref{bv:eq:cartan-K}, since
\((L+L'+L_{\mathrm{gh}})_{(1)}\) is total weight.
Every displayed identity has rational coefficients, so it holds over
the stated ring \(R\).
\end{proof}

\subsection{The strict inclusion and all coefficient changes}

The vertex subalgebra
\(Z_{R,a}=V_R^a\otimes_R G_R\) of \(K_{R,a}\) uses the vacuum in
the second matter factor.  Formula~\eqref{bv:eq:generators} shows
that it is \(Q\)-stable.  Write \(d\) for the restricted differential;
in particular,
\[
 dL=\gamma_0TL+2\gamma_1L+\frac a{12}\gamma_3.
\]
It is also the continuous Chevalley cochain complex of the Virasoro
modes \(L_n\), \(n\geq-1\), acting on \(V_R^a\): the ghost dual
to \(L_n\) is \(\gamma_{n+1}/(n+1)!\), and cochains use finite
exterior monomials.  Here continuity means finite dependence on mode
coordinates, with the vacuum module discrete.  The action sum is
finite on each vacuum state.

\begin{theorem}[Balanced vacuum comparison]\label{bv:thm:inclusion}
The inclusion
\begin{equation}\label{bv:eq:inclusion}
 I_{R,a}:Z_{R,a}\longrightarrow K_{R,a},
 \qquad v\otimes g\longmapsto v\otimes\unit\otimes g,
\end{equation}
is a strict unital dg vertex quasi-isomorphism.  Both complexes have
cohomology
\[
 R\unit\oplus R\vartheta,
 \qquad |\vartheta|=3,\qquad \vartheta^2=0,\qquad
 \vartheta=[\gamma_0\gamma_1\gamma_2].
\]
The comparison and this identification are natural under every map
of commutative rational coefficient rings, including nonflat maps.
\end{theorem}

\begin{proof}
The tensor inclusion preserves the vacuum, translation and every mode.
Its compatibility with the differential follows on the generators
from \eqref{bv:eq:generators}, and hence on all states.

Let \(\iota_0\) be left exterior differentiation with respect to
\(\gamma_1\), extended by zero on the matter states.  On the image
of \(I_{R,a}\), the free ghost contractions identify \(b_{(1)}\)
with \(\iota_0\).  Indeed
\([b_\lambda T^r\gamma]=\lambda^r\unit\), and \(b\) has no
contraction with matter.  Restricting \eqref{bv:eq:cartan-K} therefore
gives
\[
 d\iota_0+\iota_0d=H_Z,
\]
where \(H_Z\) is total weight on \(Z_{R,a}\).  On \(G_R\) this can
also be checked directly: the two terms containing \(\gamma_1\)
in \(T^r(\gamma_0\gamma_1)\) give
\((d\iota_0+\iota_0d)\gamma_r=(r-1)\gamma_r\).
The operator \(\iota_0\) is used as a chain homotopy; it is not
asserted to commute with translation.

For every nonzero integral weight \(w\), the operators
\(\iota_0/w\) and \(b_{(1)}/w\) contract the respective weight
complexes.  They are defined on every algebraic state without an
infinite sum.  The only negative-weight generator is \(\gamma_0\),
which occurs at most once.  Matter creators and \(b\)-creators have
weight at least two.  Consequently both weight-zero complexes are
exactly
\begin{equation}\label{bv:eq:four-states}
 R\{\unit,\gamma_1,\gamma_0\gamma_2,
                         \gamma_0\gamma_1\gamma_2\},
 \qquad d\gamma_1=\gamma_0\gamma_2,
\end{equation}
with every other differential zero.  The inclusion is the identity
on these four states.  This proves the cohomology assertion.  The top
state has square zero because ghost products have no mutual
contractions and are exterior products.

The contractions are \(R\)-linear.  In weight zero the additional
homotopy sends \(\gamma_0\gamma_2\) to \(\gamma_1\) and kills the
other basis states.  Thus both complexes split onto
\(R\unit\oplus R\vartheta\) by formulas preserved by any coefficient
map.  PBW identifies their scalar extensions with the same universal
constructions over the new ring.  This proves base change directly,
without requiring tensoring to preserve arbitrary quasi-isomorphisms.
\end{proof}

Taking \(R=\mathbb Q[t]\) and \(a=t\) gives a single family covering
all charges.  Its contraction uses no inverse Gram determinant.
The algebraic weight sum is essential to the stated carrier; no
completion by arbitrary infinite state sums has been introduced.

\subsection{The ghost quotient and the cubic normalization}

The inclusion \(G_R\to Z_{R,a}\) is a strict dg vertex map.  The
same contraction and four-state calculation make it a
quasi-isomorphism.  To obtain a strict map from a constant vertex
algebra, first take the quotient
\[
 E_R=G_R/(\gamma_r:r\geq3).
\]
The ideal is stable under \(T\) and \(d\).  At the only possible
low-jet exception,
\(d\gamma_3=2\gamma_1\gamma_3+\gamma_0\gamma_4\), since
\(\gamma_2^2=0\).  For \(r\geq4\), every term of
\(T^r(\gamma_0\gamma_1)\) contains a jet of order at least three.
Explicitly,
\begin{equation}\label{bv:eq:three-jets}
 \begin{gathered}
 E_R=\Lambda_R(\gamma_0,\gamma_1,\gamma_2),\qquad
 d\gamma_0=\gamma_0\gamma_1,\quad
 d\gamma_1=\gamma_0\gamma_2,\quad
 d\gamma_2=\gamma_1\gamma_2,\\
 T\gamma_0=\gamma_1,\qquad T\gamma_1=\gamma_2,\qquad T\gamma_2=0.
 \end{gathered}
\end{equation}
Its differential from degree one to degree two is an isomorphism;
the differentials out of degrees zero and two vanish.

\begin{corollary}[A strict vertex zigzag]\label{bv:cor:zigzag}
Let \(\Lambda_R(\vartheta_3)\) have zero differential and zero
translation.  There is a natural zigzag of strict unital dg vertex
quasi-isomorphisms
\begin{equation}\label{bv:eq:zigzag}
 K_{R,a}\xleftarrow{I_{R,a}}Z_{R,a}
 \xleftarrow{}G_R\xrightarrow{}E_R
 \xleftarrow{\ \vartheta_3\mapsto\gamma_0\gamma_1\gamma_2\ }
 \Lambda_R(\vartheta_3).
\end{equation}
\end{corollary}

\begin{proof}
The quotient \(G_R\to E_R\) preserves the two surviving
weight-zero classes and is a quasi-isomorphism.  In \(E_R\) the top
state is killed by both \(d\) and \(T\), and has square zero.
Its assignment from the constant exterior algebra therefore preserves
all vertex operations.  It induces both cohomology generators.
The other arrows were proved above.
\end{proof}

The quotient is needed for this last strict arrow.  In \(G_R\), and
hence in its image in \(K_{R,a}\),
\begin{equation}\label{bv:eq:translation-defect}
 T(\gamma_0\gamma_1\gamma_2)
 =\gamma_0\gamma_1\gamma_3
 =d(-\gamma_0\gamma_3)\ne0.
\end{equation}
The cohomology class is translation-invariant, but this representative
cannot be the image of a generator with zero translation under a
vertex map to the unquotiented algebra.

The factorial normalization also matters.  With
\[
 c_{-1}=\gamma_0,\qquad c_0=\gamma_1,\qquad c_1=\gamma_2/2,
 \qquad \omega_3=[c_{-1}c_0c_1],
\]
one has \(\vartheta=2\omega_3\).  To compare with the exact basis in
Proposition~\ref{prop:sl2}, set
\[
 e=L_{-1},\qquad f=-L_1,\qquad h=2L_0.
\]
The Virasoro relation gives its three \(\mathfrak{sl}_2\) brackets,
and the dual forms are
\(e^\vee=c_{-1}\), \(f^\vee=-c_1\), \(h^\vee=c_0/2\).
Their differentials agree with that proposition.  Thus its generator
has the precise image
\begin{equation}\label{bv:eq:normalization}
 c_3=[e^\vee f^\vee h^\vee]
 =\frac12\omega_3=\frac14\vartheta.
\end{equation}
This identifies the ordinary three-mode cochain algebra with
\(E_R\) after the displayed extension of coefficients and choice of
translation.  It supplies no map from the ordinary enveloping-algebra
bar to a Virasoro chiral bar.

\subsection{Relative cohomology and the charge condition}

The relative complexes impose weight zero and the kernel of
\(\iota_0\) on \(Z_{R,a}\), or of \(b_{(1)}\) on \(K_{R,a}\).
The Cartan identities make these kernels subcomplexes at weight zero.
The inclusion identifies them exactly:
\begin{equation}\label{bv:eq:relative}
 R\unit\oplus R\omega,\qquad
 \omega=\gamma_0\gamma_2,\qquad |\omega|=2,\qquad
 d=0,\quad \omega^2=0.
\end{equation}
Indeed left exterior differentiation kills these two states and sends
\(\gamma_1\) to \(\unit\) and
\(\gamma_0\gamma_1\gamma_2\) to \(-\omega\).  In the absolute
complex \(\omega=d\gamma_1\), while the cubic representative is
not relative.  The product in \eqref{bv:eq:relative} is the
zero-weight exterior-state product.  It does not make this complex a
vertex subalgebra, since
\[
 T\omega=\gamma_1\gamma_2+\gamma_0\gamma_3
\]
is a nonzero state of weight one.

Charge balance is equally exact.  Replacing the second charge by
\(26-a+\epsilon\) in \eqref{bv:eq:charge} changes the same Wick
calculation to
\[
 Q^2b=\frac\epsilon{12}\gamma_3.
\]
For every nonzero \(\epsilon\in R\), even a nonzero nilpotent,
this state is nonzero: \(R\gamma_3\) is a free PBW summand.
The four-state differential is another essential datum.  Replacing
it by zero would create false absolute classes in degrees one and
two, in addition to the unit and cubic class.

\subsection{The geometric carrier and every word length}

Now let \(k\) be a field of characteristic zero and
\(X=\mathbb A^1_k\) with its fixed coordinate.  Use the derived
category of right \(D\)-modules on the nonempty Ran space.  Its chiral
tensor is formed from external products by restriction to disjoint
configurations, direct image from that open set, and union of the
configurations.  More explicitly, for a partition
\(I=\bigsqcup_{j\in J}I_j\), let
\(j_I:U_I\hookrightarrow X^I\) exclude collisions between different
blocks.  The tensor operation on these components is
\[
 (M_j)_{j\in J}\longmapsto
 (j_I)_*(j_I)^*\left(\boxtimes_{j\in J}M_j\right).
\]
Passage to Ran uses the union maps and the compatible diagonal
pullbacks.  All these functors are derived.  We use this ambient
symmetric monoidal category, its colimit-preserving tensor and exact
filtered colimits as geometric inputs
\cite{FG12}*{Sections 1.4 and 2.1--2.4}.

For a dg vertex algebra \(V\), write \(\mathfrak a(V)\) for its
coordinate chiral Lie object, supported on the singleton diagonal in
Ran.  On \(X\) its underlying right \(D\)-module has
\(\mathcal O_X\)-module \(V\otimes_k\omega_X\), with connection
determined by translation and chiral operation determined by the full
vertex operation.  This is the coordinate vertex--chiral
correspondence \cite{BD04}*{Section 0.15}.  A strict dg vertex map
induces a chiral map.  Its underlying complex construction is exact
over \(k\), so a vertex quasi-isomorphism induces a chiral
quasi-isomorphism.  No derived equivalence of arbitrary module
categories is needed here.

For coaugmentation, formally adjoin the empty-configuration sector;
denote its tensor unit by \(\unit_{\varnothing}\).  The vertex vacuum
remains an element of \(\mathfrak a(V)\), supported on the singleton
diagonal, and is distinct from this empty-sector unit.

Let \(A\) be a dg Lie object for the chiral tensor.  Define its
direct-sum Chevalley object by
\begin{equation}\label{bv:eq:CE}
 \CE_{\oplus}^{\mathrm{ch}}(A)
 =\unit_{\varnothing}\oplus
   \bigoplus_{r\geq1}\Sym_{\mathrm{ch}}^r(A[1]).
\end{equation}
Tensor products and finite symmetric-group coinvariants are derived.
For a homogeneous element \(a\), write \(sa\) for its copy in
\(A[1]\), of degree \(|a|-1\).  The internal differential and the
suspended bracket have corestrictions
\[
 q_1(sa)=-s(da),\qquad
 q_2(sa\,sb)=(-1)^{|a|}s[a,b].
\]
Extend \(q_1\) to each word and sum \(q_2\) over all unordered
pairs, moving that pair together with its Koszul sign.  The resulting
differential preserves or lowers word length.  The coproduct is the
sum over all unshuffles, including the empty subword.

The differential squares to zero: the chain identity for the bracket
cancels the mixed terms; disjoint pair contractions cancel in the two
orders; and contractions sharing an input sum to the shifted Jacobi
identity.  Performing an unshuffle before or after a pair contraction
gives the same terms, so the differential is a coderivation.  Each
word has finitely many inputs, and its sufficiently iterated reduced
coproduct is zero.  Thus \eqref{bv:eq:CE} is a conilpotent chiral
coalgebra.

\begin{proposition}[Direct-sum Chevalley invariance]
\label{bv:prop:CE}
A quasi-isomorphism \(f:A\to B\) of these dg Lie objects induces a
quasi-isomorphism of coaugmented conilpotent chiral coalgebras
\[
 \CE_{\oplus}^{\mathrm{ch}}(f):
 \CE_{\oplus}^{\mathrm{ch}}(A)\longrightarrow
 \CE_{\oplus}^{\mathrm{ch}}(B).
\]
There is no boundedness assumption on their cohomological degrees.
\end{proposition}

\begin{proof}
Filter by the finite sums
\[
 F_r\CE_{\oplus}^{\mathrm{ch}}(A)
 =\bigoplus_{0\leq j\leq r}\Sym_{\mathrm{ch}}^j(A[1]),
 \qquad \Sym_{\mathrm{ch}}^0(A[1])=\unit_{\varnothing}.
\]
They are subcomplexes, since the bracket lowers length.  The quotient
\(F_r/F_{r-1}\) is \(\Sym_{\mathrm{ch}}^r(A[1])\) with its internal
differential.  Derived tensors and derived finite-group coinvariants
preserve equivalences; hence the map on each quotient is an
equivalence.  Starting with the identity on \(F_0\), the exact
triangle for two consecutive stages proves the assertion for every
\(F_r\).  Exact filtered colimit gives it for the direct sum.
Naturality preserves every unshuffle and the bracket differential,
which proves the coalgebra assertion.
\end{proof}

Apply \(\mathfrak a\) and then \(\CE_{\oplus}^{\mathrm{ch}}\) to
each arrow of \eqref{bv:eq:zigzag}, with \(R=k\).  The resulting
zigzag compares the balanced vacuum with
\(\Lambda_k(\vartheta_3)\) as conilpotent chiral coalgebras, in every
word length and every ghost degree.  The proof uses finite increasing
stages and their exact union, rather than convergence of an unbounded
spectral sequence.

At the constant exterior endpoint the singular Lie conformal bracket
vanishes, but the geometric chiral operation retains the regular
product.  In its coordinate residue description, a simple pole with
coefficient \(u\otimes v\) gives the product \(uv\) on the diagonal;
the vacuum product is already nonzero.  Discarding these products
would change the chiral object
\cite{BD04}*{Sections 3.3.2--3.3.3}.

The direction of the filtration also fixes the scope.  For a nonzero
bracket, a word of length \(r+1\) can have a differential component
of length \(r\).  The naive projection \(F_{r+1}\to F_r\) is then
not a chain map.  A completed product requires a different compatible
filtration and its own convergence argument.  Likewise, the chiral
tensor in \eqref{bv:eq:CE} is not the diagonal tensor used for the
ordered associative bar.  An identification with that bar requires
an additional comparison of the tensors and their operations.

Finally, the first matter stress tensor satisfies
\(QL=\gamma TL+2(T\gamma)L+(a/12)T^3\gamma\ne0\).
Thus the matter embedding from \(V_k^a\) with zero differential is
not a dg map.  The unmarked zigzag does not supply a marked boundary
resolution, a compatible action of a transverse monoid, or a finality
property for such actions.  Its positive conclusion is the explicit
vertex and Lie--chiral comparison above; the ordered collision action
and the root-graded comparison retain their separate construction
problems.
